\begin{center}
\begin{tikzpicture}[scale=0.8, >=latex]
    \draw[very thin, color=gray!30, xstep=pi/4, ystep=1] (-1.7,-3.5) grid (6.5,4.5);
    \draw[->] (-1.8,0) -- (7.0,0) node[right] {$x$};
    \draw[->] (0,-3.5) -- (0,4.5) node[above] {$y$};
    \foreach \x/\xtext in {-1.57/-\frac{\pi}{2}, 1.57/\frac{\pi}{2}, 3.14/\pi, 4.71/\frac{3\pi}{2}, 6.28/2\pi}
        \draw (\x,2pt) -- (\x,-2pt) node[below, font=\small, fill=white, inner sep=1pt] {$\xtext$};
    \foreach \y in {-3,-2,-1,1,2,3,4}
        \draw (2pt,\y) -- (-2pt,\y) node[left, font=\small, fill=white, inner sep=1pt] {\y};
    \node[below left, font=\small] at (0,0) {$0$};
    % f(x) = 3*sin(2x)
    \draw[very thick, color=maincolor, samples=200, domain=-1.6:6.3] plot (\x, {3*sin(2*deg(\x))}) node[above right] {$f$};
\end{tikzpicture}
\end{center}
\Askhsh{15062}{4}
\begin{ylh}{1}
\item 2.1 Μονοτονία-Ακρότατα-Συμμετρίες Συνάρτησης
\item 3.4 Οι τριγωνομετρικές συναρτήσεις
\end{ylh}
Στο διπλανό σχήμα δίνεται η γραφική παράσταση μιας συνάρτησης f που είναι της μορφής
$f(x) = \rho\hm (\alpha x), x \in \mathbb{R}$ και $\alpha, \rho > 0$
\begin{alist}
\item Να βρείτε, με βάση το σχήμα, την περίοδό της, την μέγιστη και την ελάχιστη τιμή της.
\monades{6}
\item Με βάση τις απαντήσεις στο προηγούμενο ερώτημα, να βρείτε τους αριθμούς α και ρ.
\monades{6}
Έστω $\rho = 3$ και $\alpha = 2$. Θεωρούμε επίσης τη συνάρτηση $g(x) = x^4 - 2x^2 + 5, x \in \mathbb{R}$.
\item Να αποδείξετε ότι η ελάχιστη τιμή της είναι ίση με 4.
\monades{7}
\item Να αιτιολογήσετε γιατί οι γραφικές παραστάσεις των f, g δεν έχουν κοινό σημείο.
\monades{6}
\end{alist}

